\section{Data Sourcing and Processing}
\label{sec:data}

% Rubric (Problem & Dataset): "Dataset described with sources, limits, and relevance."
% Guidelines:
%   - Where the data comes from.
%   - Preprocessing (missing data, discretization, feature selection) and WHY
%     each step matters for the problem.
%
% Outline:
% - Source and description \citep{vito2008airquality, devito2008calibration}:
%   size, time span, variables (table of variables used).
% - Missing data: -200 encoding; missingness per variable (table or figure);
%   how handled (EM vs. drop) and why.
% - Discretization: bin scheme (quantile / domain thresholds), number of bins,
%   sensitivity checked in Sec. 4.
% - Feature selection: which sensors / pollutants kept (e.g., drop NMHC) and why.
% - Train / test split: time-based (no shuffling of hourly data) and why.
% - Limitations: one device, one site, one year; sensor drift.

\todo{Data source, preprocessing steps, and justification.}

% Example table:
% \begin{table}[t]
%   \centering
%   \caption{Variables used.}
%   \label{tab:variables}
%   \begin{tabular}{lll}
%     \toprule
%     Variable & Type & Role in model \\
%     \midrule
%     CO(GT) & reference analyzer & partially observed \\
%     PT08.S1 & tin-oxide sensor & observed \\
%     \bottomrule
%   \end{tabular}
% \end{table}
